\sect{फाल्गुन-शुक्लामलक्येकादशी-माहात्म्यम्}
\label{sec:padma-phalguna-shuklamalaki}

\ganapatyadiDhyanam
\padmaPuranam

\dnsub{अथ सप्तचत्वारिंशोऽध्यायः॥४७}

\uvacha{श्रीकृष्ण उवाच}

\twolineshloka
{माहात्म्यं विजयायाश्च श्रुतं कृष्ण महत्फलम्}
{फाल्गुनस्यार्जुने पक्षे यन्नाम्नी तां वदाधुना}% ॥१॥

\uvacha{श्रीकृष्ण उवाच}

\twolineshloka
{धर्मपुत्र महाभाग शृणु वक्ष्यामि तेऽधुना}
{योक्ता पृष्टेन मान्धात्रा वसिष्ठेन महात्मना}% ॥२॥

\twolineshloka
{फाल्गुनस्य विशेषेण विशेषः कथितो नृप}
{आमलकीव्रतं पुण्यं विष्णुलोकफलप्रदम्}% ॥३॥

\twolineshloka
{आमलक्यास्तले गत्वा जागरं तत्र कारयेत्}
{कृत्वा जागरणं रात्रौ गोसहस्रफलं लभेत्}% ॥४॥

\uvacha{मान्धातोवाच}

\twolineshloka
{आमलकी कदा ह्येषा उत्पन्ना द्विजसत्तम}
{एतत् सर्वं ममाऽऽचक्ष्व परं कौतूहलं हि मे}% ॥५॥

\twolineshloka
{कस्मादियं पवित्रा च कस्मात् पापप्रणाशिनी}
{कस्माज्जागरणं कृत्वा गोसहस्रफलं लभेत्}% ॥६॥

\uvacha{वसिष्ठ उवाच}

\twolineshloka
{कथयामि महाभाग यथेयमभवत्क्षितौ}
{आमलकी महावृक्षः सर्वपापप्रणाशनः}% ॥७॥

\twolineshloka
{एकार्णवे पुरा जाते नष्टे स्थावरजङ्गमे}
{नष्टे देवासुरगणे प्रणष्टोरगराक्षसे}% ॥८॥

\twolineshloka
{तत्र देवादिदेवेशः परमात्मा सनातनः}
{जगाम ब्रह्मपरममात्मनः पदमव्ययम्}% ॥९॥

\twolineshloka
{ततोऽस्य जाग्रतो ब्रह्ममुखाच्छशिसमप्रभः}
{ष्ठीवनाद् बिन्दुरुत्पन्नः स भूमौ निपपात ह}% ॥१०॥

\twolineshloka
{तस्माद् बिन्दोः समुत्पन्नः स्वयं धात्री नगो महान्}
{शाखाप्रशाखाबहुलः फलभारेण नामितः}% ॥११॥

\twolineshloka
{सर्वेषां चैव वृक्षाणामादिरोहः प्रकीर्तितः}
{ब्रह्मणाऽथ ततः पश्चात् संसृष्टाश्च इमाः प्रजाः}% ॥१२॥

\twolineshloka
{देवदानवगन्धर्वयक्षराक्षसपन्नगान्}
{असृजद् भगवान् देवो महर्षींश्च तथाऽमलान्}% ॥१३॥

\twolineshloka
{आजग्मुस्तत्र देवास्ते यत्र धात्री हरिप्रिया}
{तां दृष्ट्वा ते महाभाग परं विस्मयमागताः}% ॥१४॥

\twolineshloka
{न जानीम इमं वृक्षं चिन्तयन्तोऽभिसंस्थिताः}
{एवं चिन्तयतां तेषां वागुवाचाशरीरिणी}% ॥१५॥

\twolineshloka
{आमलकी नगो ह्येष प्रवरो वैष्णवो मतः}
{अस्य संस्मरणादेव लभेद् गोदानजं फलम्}% ॥१६॥

\twolineshloka
{स्पर्शनाद् द्विगुणं पुण्यं त्रिगुणं धारणात् तथा}
{तस्मात् सर्वप्रयत्नेन सेव्या आमलकी सदा}% ॥१७॥

\twolineshloka
{सर्वपापहरा प्रोक्ता वैष्णवी पापनाशिनी}
{तस्या मूले स्थितो विष्णुस्तदूर्ध्वे च पितामहः}% ॥१८॥

\twolineshloka
{स्कन्धे च भगवान् रुद्रः संस्थितः परमेश्वरः}
{शाखासु मुनयः सर्वे प्रशाखासु च देवताः}% ॥१९॥

\twolineshloka
{पर्णेषु चाऽऽसते देवाः पुष्पेषु मरुतस्तथा}
{प्रजानां पतयः सर्वे फलेष्वेव व्यवस्थिताः}% ॥२०॥

\twolineshloka
{सर्वदेवमयी ह्येषा धात्री च कथिता मया}
{तस्मात् पूज्यतमा ह्येषा विष्णुभक्तिपरायणैः}% ॥२१॥

\uvacha{ऋषय ऊचुः}

\twolineshloka
{को भवान्न हि जानीमः कस्मात् कारणतां गतः}
{देवो वा यदि वा चान्यः कथयस्व यथातथम्}% ॥२२॥

\uvacha{वागुवाच}

\twolineshloka
{यः कर्ता सर्वभूतानां भुवनानां च सर्वशः}
{विस्मितान् विदुषः प्रेक्ष्य सोऽहं विष्णुः सनातनः}% ॥२३॥

\twolineshloka
{तच्छ्रुत्वा देवदेवस्य भाषितं ब्रह्मणः सुताः}
{अनादिनिधनं देवं स्तोतुं तत्र प्रचक्रमुः}% ॥२४॥

\twolineshloka
{नमो भूतात्मभूताय आत्मने परमात्मने}
{अच्युताय नमो नित्यमनन्ताय नमो नमः}% ॥२५॥

\twolineshloka
{दामोदराय कवये यज्ञेशाय नमो नमः}
{एवं स्तुतस्तु ऋषिभिस्तुतोष भगवान् हरिः}% ॥२६॥

\onelineshloka*
{प्रत्युवाच महर्षींस्तानभीष्टं किं ददामि वः}

\uvacha{ऋषय ऊचुः}

\onelineshloka
{यदि तुष्टोऽसि भगवन्नस्माकं हितकाम्यया}% ॥२७॥

\twolineshloka
{व्रतं किञ्चित् समाख्याहि स्वर्गमोक्षफलप्रदम्}
{धनधान्यप्रदं पुण्यमात्मनस्तुष्टिकारकम्}% ॥२८॥

\twolineshloka
{अल्पायासं बहुफलं व्रतानामुत्तमं व्रतम्}
{कृतेन येन देवेश विष्णुलोके महीयते}% ॥२९॥

\uvacha{विष्णुरुवाच}

\twolineshloka
{फाल्गुने शुक्लपक्षे तु पुष्येण द्वादशी यदि}
{भवेत् सा च महापुण्या महापातकनाशिनी}% ॥३०॥

\twolineshloka
{विशेषस्तत्र कर्तव्यः शृणुध्वं द्विजसत्तमाः}
{आमलकीं च सम्प्राप्य जागरं तत्र कारयेत्}% ॥३१॥

\threelineshloka
{सर्वपापविनिर्मुक्तो गोसहस्रफलं लभेत्}
{एतद् वः कथितं विप्रा व्रतानां व्रतमुत्तमम्}
{अर्चयित्वाऽच्युतं तस्यां विष्णुलोकान्न मुच्यते}% ॥३२॥

\uvacha{ऋषय ऊचुः}

\twolineshloka
{व्रतस्यास्य विधिं ब्रूहि परिपूर्णं कथं भवेत्}
{के मन्त्राः के नमस्काराः देवता का प्रकीर्तिता}% ॥३३॥

\twolineshloka
{कथं दानं कथं स्नानं कश्च पूजाविधिः स्मृतः}
{अर्घार्चनस्य मन्त्रं तु कथयस्व यथातथम्}% ॥३४॥

\uvacha{विष्णुरुवाच}

\twolineshloka
{श्रूयतां यो विधिः सम्यग् व्रतस्यास्य द्विजर्षभाः}
{एकादश्यां निराहारः स्थित्वा चैव परेऽहनि}% ॥३५॥

\twolineshloka
{भोक्ष्येऽहं पुण्डरीकाक्ष शरणं मे भवाच्युत}
{इति कृत्वा तु नियमं दन्तधावनपूर्वकम्}% ॥३६॥

\twolineshloka
{नाऽऽलपेत् पतितांश्चौरांस्तथा पाषण्डिनो नरान्}
{दुर्वृत्तान् भिन्नमर्यादान् गुरुदारप्रधर्षकान्}% ॥३७॥

\twolineshloka
{अपराह्णे ततः स्नानं विधिना कारयेद् बुधः}
{नद्यां तडागे कूपे वा गृहे वा नियतात्मवान्}% ॥३८॥

\twolineshloka
{मृत्तिकालम्भनं पूर्वं ततः स्नानं च कारयेत्}
{अश्वक्रान्ते रथक्रान्ते विष्णुक्रान्ते वसुन्धरे}% ॥३९॥

\onelineshloka
{मृत्तिके हर मे पापं यन्मया दुष्कृतं कृतम्}% ॥४०॥
[इति मृत्तिकामन्त्रः]

\twolineshloka
{त्वमम्बु सर्वभूतानां जीवनं तनुरक्षकम्}
{स्वेदजोद्भिजजातीनां रसानां पतये नमः}% ॥४१॥

\twolineshloka
{स्नातोऽहं सर्वतीर्थेषु ह्रदप्रस्रवणेषु च}
{नदीषु देवखातेषु इदं स्नानं तु मे भवेत्}% ॥४२॥
[इति स्नानमन्त्रः]

\twolineshloka
{जामदग्न्यं मुनिं चैव कारयित्वा हिरण्मयम्}
{मासकस्य सुवर्णस्य तदर्धार्धेन वा पुनः}% ॥४३॥

\twolineshloka
{गृहमागत्य पूजायाः पूजाहोमं तु कारयेत्}
{ततश्चाऽऽमलकीं गच्छेत् सर्वोपस्करसंयुतः}% ॥४४॥

\twolineshloka
{आमलकीं ततो गत्वा परिशोध्य समन्ततः}
{स्थापयेत् सततं कुम्भमव्रणं मन्त्रपूर्वकम्}% ॥४५॥

\twolineshloka
{पञ्चरत्नसमोपेतं दिव्यगन्धाधिवासितम्}
{छत्रोपानद्युगोपेतं सितचन्दनचर्चितम्}% ॥४६॥

\twolineshloka
{स्रग्दामलम्बितग्रीवं सर्वधूपविधूपितम्}
{दीपमालाकुलं कुर्यात् सर्वतः सुमनोहरम्}% ॥४७॥

\twolineshloka
{तस्योपरि न्यसेत् पात्रं दिव्यलाजैः प्रपूरितम्}
{पात्रोपरि न्यसेद् देवं जामदग्न्यं महाप्रभम्}% ॥४८॥

\twolineshloka
{विशोकाय नमः पादौ जानुनी विश्वरूपिणे}
{उग्राय च ततोऽप्यूरू कटी दामोदराय च}% ॥४९॥

\twolineshloka
{उदरं पद्मनाभाय उरः श्रीवत्सधारिणे}
{चक्रिणे वामबाहुं च दक्षिणं गदिने नमः}% ॥५०॥

\twolineshloka
{वैकुण्ठाय नमः कण्ठमास्यं यज्ञमुखाय वै}
{नासां विशोकनिधये वासुदेवाय चाक्षिणी}% ॥५१॥

\twolineshloka
{ललाटं वामनायेति रामायेति भ्रुवौ नमः}
{सर्वात्मने तु तच्छीर्षं नम इत्यभिपूजयेत्}% ॥५२॥
[इति पूजामन्त्रः]

\twolineshloka
{ततो देवाधिदेवाय अर्घ्यं चैव प्रदापयेत्}
{फलेन चैव शुभ्रेण भक्तियुक्तेन चेतसा}% ॥५३॥

\twolineshloka
{नमस्ते देवेदेवेश जामदग्न्य नमोऽस्तु ते}
{गृहाणार्घ्यमिमं दत्तमामलक्या युतो हरेः}

[इत्यर्घ्यमन्त्रः]

\twolineshloka
{ततो जागरणं कुर्याद् भक्तियुक्तेन चेतसा}
{नृत्यैर्गीतैश्च वादित्रैर्धर्माख्यानैः स्तवैरपि}% ॥५४॥

\twolineshloka
{वैष्णवैश्च तथाऽऽख्यानैः क्षपयेत् सर्वशर्वरीम्}
{प्रदक्षिणां ततः कुर्याद् धात्र्या वै विष्णुनामभिः}% ॥५५॥

\twolineshloka
{अष्टाधिकं शतं चैव अष्टाविंशतिरेव वा}
{ततः प्रभाते समये कृत्वा नीराजनं हरेः}% ॥५६॥

\twolineshloka
{ब्राह्मणं पूजयित्वा तु सर्वं तस्मै निवेदयेत्}
{जामदग्न्यं घटं तत्र वस्त्रयुग्ममुपानहौ}% ॥५७॥

\twolineshloka
{जामदग्न्यस्वरूपेण प्रीयतां मम केशवः}
{ततश्चाऽऽमलकीं स्पृष्ट्वा कृत्वा चैव प्रदक्षिणाम्}% ॥५८॥

\twolineshloka
{स्नानं कृत्वा विधानेन ब्राह्मणान् भोजयेत् ततः}
{ततश्च स्वयमश्नीयात् कुटुम्बेन समावृतः}% ॥५९॥

\twolineshloka
{एवं कृतेन यत् पुण्यं तत् सर्वं कथयामि ते}
{सर्वतीर्थेषु यत् पुण्यं सर्वदानेषु यत्फलम्}% ॥६०॥

\twolineshloka
{सर्वयज्ञाधिकं चैव लभते नात्र संशयः}
{एतद् वः सर्वमाख्यातं व्रतानामुत्तमं व्रतम्}% ॥६१॥

\twolineshloka
{एतावदुक्त्वा देवेशस्तत्रैवान्तरधीयत}
{ते चापि ऋषयः सर्वे चक्रुः सर्वमशेषतः}% ॥६२॥

\twolineshloka
{तथा त्वमपि राजेन्द्र कर्तुमर्हसि सत्तम}
{व्रतमेतद्दुराधर्षं सर्वपापप्रमोचनम्}% ॥६३॥

॥इति श्रीपाद्मे महापुराणे पञ्चपञ्चाशत्साहस्र्यां संहितायामुत्तरखण्डे उमापतिनारदसंवादान्तर्गतकृष्णयुधिष्ठिरसंवादे फाल्गुन-शुक्लामलक्येकादशी-माहात्म्यं नाम सप्तचत्वारिंशोऽध्यायः॥४७॥

\genMangalaShloka

\hyperref[sec:ekadashi_mahatmyam_padma_puranam]{\closesub}
\clearpage

